\chapter{积分学}
\section{这样评论黎曼公正吗？}

\noindent\textbf{1. 不连续函数进入了人们的视野}\quad
迪厄多内（Jean Dieudonné）在《现代分析基础》（Éléments d'Analyse）第一卷（中译本由科学出版社1982年出版）159页说了这样一段话，以解释为什么这部大书（全书共九卷）中完全不讲黎曼积分：“如果不是它的有权威的名字，它老早就该没落下去了。因为对任何一位从事研究工作的数学家来说（带着对黎曼的天才应有的尊敬），十分清楚，现今这一‘理论’的重要性在测度与积分的一般理论中，最多不过是一普通的有趣的练习（参看13.9问题7），只有那种学究传统的顽固的保守主义才会把它冻结成课程的正规部分，长时间以后必将失去它的历史重要性。”现在我们要追溯一下历史，看一看黎曼是在什么样的历史背景下提出他的积分理论，以及他的理论的不足之处究竟在什么地方。

这就需要从傅里叶讲起。傅里叶从研究大量物理问题（主要是热传导问题）出发，指出“任何函数”$f(x),x\in[-\pi,\pi]$都可以写成一个无穷级数
\begin{equation}
 f(x)=\frac{a_0}{2}+\sum_{n=1}^{\infty}\bigl(a_n\cos nx+b_n\sin nx\bigr),
 \tag{1}
\end{equation}
这里
\begin{equation}
 a_n=\frac1\pi\int_{-\pi}^{\pi}f(\xi)\cos n\xi\,d\xi,
 \qquad
 b_n=\frac1\pi\int_{-\pi}^{\pi}f(\xi)\sin n\xi\,d\xi.
 \tag{2}
\end{equation}

在傅里叶（1768--1830）的时代（19世纪初），什么是“任何函数”是一个没有说清的问题，本书第二章讲到了这件事。当时，一种说法是：函数就是一个解析表达式。因此，不同的解析表达式就表示不同的函数；另一种说法则是：顺手一笔画出来的曲线就是一个函数。这是同一个人欧拉，在不同场合下讲的话。傅里叶在他的名著《热的解析理论》（Jean-Baptiste Joseph Fourier, \textit{Théorie Analytique de la Chaleur}，中译本武汉出版社1993年出版）中讨论三个函数的一般理论时（中译本132--133页），就给出了一个例子（见图4-1-1）
\begin{equation}
 \cos x-\frac13\cos3x+\frac15\cos5x-\cdots=
 \begin{cases}
 \dfrac{\pi}{4},&|x|<\dfrac{\pi}{2},\\[4pt]
 -\dfrac{\pi}{4},&\dfrac{\pi}{2}<|x|<\pi.
 \end{cases}
 \tag{3}
\end{equation}

\begin{figure}[htbp]
\centering
\begin{tikzpicture}[x=1.45cm,y=3.3cm,bella figure]
  \draw[->] (-3.7,0)--(3.7,0);
  \draw[->] (0,-.52)--(0,.52);
  \draw[thick] (-3.5,-.25)--(-1.57,-.25);
  \draw[thick] (-1.57,.25)--(1.57,.25);
  \draw[thick] (1.57,-.25)--(3.5,-.25);
  \draw[dashed,thick] (-1.57,-.25)--(-1.57,.25);
  \draw[dashed,thick] (1.57,.25)--(1.57,-.25);
  \node[below] at (-3.1416,0) {$-\pi$};
  \node[below] at (-1.5708,0) {$-\frac\pi2$};
  \node[below right] at (0,0) {$O$};
  \node[below] at (1.5708,0) {$\frac\pi2$};
  \node[left] at (0,.25) {$\frac\pi4$};
  \node[left] at (0,-.25) {$-\frac\pi4$};
\end{tikzpicture}
\caption*{图4-1-1}
\end{figure}

同一个级数怎么能表示不同的函数呢？看来并不是所有的函数都如雅可比说的是“合理的”函数，而可以是不连续的，而且研究这种有不连续点的函数，在物理上也是“合理的”。这一个例子以及后来越来越多的例子才引导到关于一般函数的“狄利克雷定义”，即本书第二章所介绍的定义。越来越多的例子表明，函数的不连续点的集合构造可以是十分复杂的。例如“狄利克雷函数”
\begin{equation}
 \chi(x)=
 \begin{cases}
 1,&x\text{ 为有理数},\\
 0,&x\text{ 为无理数}
 \end{cases}
 \tag{4}
\end{equation}
是处处不连续的，而形状十分相近的“黎曼函数”
\begin{equation}
 R(x)=
 \begin{cases}
 1/q,&x=p/q\text{ 为既约有理数},\\
 0,&x\text{ 为无理数}
 \end{cases}
 \tag{5}
\end{equation}
却只在有理点不连续。其实傅里叶所说的“任何函数”是指的连续函数，但是如(3)那样的函数在傅里叶看来也是连续的，因为它的图像如图4-1-1所示是一条连续曲线，所以表示连续函数。傅里叶把图1中的虚线也包括到(3)的图像中去了，而按狄利克雷的定义这样做是不许可的，因为这样一来在例如 $x=-\pi/2$ 处，函数的值就不确定了。再用柯西关于连续函数的定义才看出这个点是不连续点。不过傅里叶的书实际上成于1810左右，狄利克雷的函数定义和柯西关于连续性的定义乃是1820--1830左右的事，何况他们二人自己有时也会“糊涂”的。不论如何，当函数的不连续性逼迫人们予以正视时，从定义傅里叶系数 $a_n$ 与 $b_n$ 的(2)式看，有必要研究不连续函数的积分。

把定积分定义为积分和的极限
\begin{equation}
 \int_a^b f(x)\,dx=\lim\sum f(\xi_i)(x_i-x_{i-1})
 \tag{6}
\end{equation}
的是柯西，但是他假设了 $f(x)$ 在 $[a,b]$ 上连续。如果 $f(x)$ 在 $x=c$ 时有间断，则定义
\[
 \int_a^b f(x)\,dx=\int_a^c f(x)\,dx+\int_c^b f(x)\,dx.
\]
现在用的反常积分的定义也是柯西给出的。但是比较一般的函数 $f(x)$ 的积分如何定义，则是黎曼1854年在研究三角级数（即形如(1)之右方的级数，但是并不确定是否有一个 $f(x)$ 存在使(2)成立）时提出来的，当时黎曼已意识到，这一项研究虽然不一定有直接的物理学上的应用，在数学上却是重要的。

上面我们看见了两个特例(4)与(5)，这使人产生一个疑问：为什么要研究这些人为的特例？上一章\S7中我们引用了庞加莱的话，然后又引用了佩兰关于布朗运动的议论。那么，研究这些古怪的函数是为了证明自己前辈的无能吗？难道黎曼（还有魏尔斯特拉斯）还需要靠“揭发”自己前辈的无能才能巩固自己的地位吗？或者是不是这些函数与某些物理现象有关？都不是。上一章讲到佩兰关于处处不可微函数与布朗运动的关系的话，当时只是一种设想而已。有一种偏见，以为数学的发展是一帆风顺的：从几个不辩自明的真理（公理、公设），以及任意设定的某些对象的性质（定义）出发，按照一些确定的推理规则，就可以得出一个又一个确定无疑的真理。从这个意义上说，数学的方法论是“先验”的。上一章\S6中我们也说过这个意思。但是数学之为“先验”只是相对于其它科学之“后验”而言。上一章我们正是相对于力学中最小作用的原理的演变讲数学方法论的先验性质的。其实，数学的发展同样充满了争论和矛盾。本书各章都说明这个问题。前辈数学家发现了一些真理，以为它是正确的；后来的数学家发现了其中有不少毛病、漏洞，而且这些毛病和漏洞常以反例形式出现。这样，人们去清理前辈的工作，补足他们的缺失，发展他们的成就。甚至许多定义也不是确定于证明之前，而是产生于证明过程之中。举一个例，早在莱布尼茨就认为如果连续函数的级数 $\sum_{n=1}^{\infty}u_n(x)$ 收敛，则其和也一定连续。柯西也是这样想的，而且（怀着对莱布尼茨天才应有的尊敬）力图去证明它，甚至傅里叶的例子(3)开始也不被认为是反例，因为傅里叶对和函数的图形补上了一些虚线，使之成为“连续曲线”。直到柯西给出连续函数的定义后，才知道这个连续曲线并非连续函数的图像。但是柯西仍然力图去证明那个可以追溯到莱布尼茨的“定理”，他的这种努力甚至见于他为数学分析严格性奠基的力作《代数分析教程》（\textit{Cours d'Analyse Algébrique}, 1821）。直到1826年，阿贝尔指出，柯西的“定理”有例外（定理怎么能有例外？）！\jiaokan{原文作“1876年”。阿贝尔于1829年逝世；这里所指的是他1826年对柯西连续函数级数定理提出“例外”的工作，故改为1826年。}最终打破这个哑谜的是赛德尔（P. L. Seidel），他在1847年分析柯西用 $\varepsilon-N$ 语言来表达的证明时指出，这个 $N$ 不仅依赖于 $\varepsilon$，还依赖于 $x$：$N=N(\varepsilon,x)$。如果 $x$ 只能取有限多个值，从这些 $N(\varepsilon,x)$ 中自然可以挑出最大的一个作 $N(\varepsilon)$。但若 $x$ 可以取无穷多个值（例如某一区间 $[a,b]$ 中的一切值），则 $N(\varepsilon,x)$ 就不一定有有限的最大值 $N(\varepsilon)$ 以适合一切 $x$ 点了，因此柯西的证明失效。赛德尔在自己的论文中说了一段话：“方才已经确认，定理不是普遍有效的，因为它的证明必定依赖某个额外的隐蔽假定（即 $N(\varepsilon,x)$ 对一切 $x$ 有有限的最大值——本书作者）。有鉴于此，我们要给证明作一次更细微的分析，发现那个隐藏的假设倒不很难，于是可以反推这个假设所表达的条件确是表示不连续函数的级数满足不了的。因为，只有这样，那串别无他错的证明才跟另一头的确凿事实才能重新归于一致。”这就是说，从分析柯西的证明失误人们找出了漏洞，发现了收敛的级数应分为两类，一类是 $N(\varepsilon,x)$ 有最大值的，另一类则是没有的。至此，一致收敛性的定义呼之欲出：并不是人们先验地有了一致收敛性定义，并由之证明了相应定理，而是从证明过程中这个定义自己跑到我们面前。我们只好承认它，应用它。当时看到了这一点的不只有赛德尔，还有斯托克斯、魏尔斯特拉斯稍早一些也在幂级数研究中独立地得到了 一致收敛性概念。到底是谁首先明确宣布这个定义，现在人们一般地归功于魏尔斯特拉斯。不过斯托克斯、赛德尔还有傅里叶本人都指出它与收敛速度有关。（以上材料引自拉卡托斯《证明与反驳》，附录一，中译本152--170页）。一致收敛的案例说明了数学的发展并不是先验的，反例的重要性在于它揭露了无法回避的矛盾，数学的发展并不是为了说明种种反例，相反，研究反例是为了发展数学。19世纪最后二三十年出现了大量“反例”，正表明数学的深入发展，在数学对世界的认识的过程中起了极大作用。这可以说是现代数学的一大特点。黎曼积分正是这种历史背景下出现的第一个系统的积分理论。回忆到上一章讲的处处不可微函数的例子是魏尔斯特拉斯1861年发现的，可以看到，不连续函数进入人们的视野是时代的产物。

\noindent\textbf{2. 函数的可积性}\quad
黎曼第一个提出可积性问题，即要求刻画出使(6)之极限存在的函数类。（这里说明一下，本书中讲到的 $f(x)$ 都是实值函数，若 $f(x)$ 是复值的，则划分开实虚部处理。）人们通常都说“积分是积分和的极限”，这个说法并不准确，我们学过了序列的极限与函数的极限，但(6)式右方的积分和
\begin{equation}
 \sum_{i=1}^{n} f(\xi_i)(x_i-x_{i-1})
 \tag{7}
\end{equation}
既不构成序列，也不是某连续变量的函数。为了明确地了解它，我们首先要求 $f(x)$ 定义在有限区间 $[a,b]$ 上，而且 $f(x)$ 不可以无界。对 $[a,b]$ 作一个分划 $P:a=x_0<x_1<\cdots<x_n=b$，并以 $\delta_P$ 表示 $\max(x_i-x_{i-1})$，因为 $f(x)$ 在 $[a,b]$ 上有界，从而在每个小区间 $[x_{i-1},x_i]$ 上也有界，从而有上（下）确界 $M_i(m_i)$，$-\infty<m_i\le M_i<+\infty$。反之，若 $f(x)$ 无界，$M_i,m_i$ 中必有 $\pm\infty$，这会使下式没有意义。我们不考虑 $f(x)$ 在 $[x_{i-1},x_i]$ 中是否可以达到 $m_i$ 与 $M_i$，而定义 $f(x)$ 相应于分划 $P$ 的上和 $S_P$ 与下和 $s_P$：
\begin{equation}
 S_P=\sum_{i=1}^{n}M_i(x_i-x_{i-1}),\qquad
 s_P=\sum_{i=1}^{n}m_i(x_i-x_{i-1}).
 \tag{8}
\end{equation}
很清楚
\begin{equation}
 s_P\le S_P.
 \tag{9}
\end{equation}
但对不同的分划 $P'$ 与 $P''$，$s_{P'}$ 与 $S_{P''}$ 可否比较？首先我们注意到，若在分划 $P$ 中加一个分点 $x'$（称为 $P$ 的“加细”，加细后的新分划为 $P'$）例如，$x_{i-1}<x'<x_i$，则 $s_P$ 和 $S_P$ 变成
\begin{align*}
 s_{P'}={}&\sum_{j=1}^{i-1}m_j(x_j-x_{j-1})+m_i'(x'-x_{i-1})+m_i''(x_i-x')
 +\sum_{j=i+1}^{n}m_j(x_j-x_{j-1}),\\
 S_{P'}={}&\sum_{j=1}^{i-1}M_j(x_j-x_{j-1})+M_i'(x'-x_{i-1})+M_i''(x_i-x')
 +\sum_{j=i+1}^{n}M_j(x_j-x_{j-1}).
\end{align*}
这里
\begin{align*}
 m_i'&=\inf_{[x_{i-1},x']}f(x), & m_i''&=\inf_{[x',x_i]}f(x),\\
 M_i'&=\sup_{[x_{i-1},x']}f(x), & M_i''&=\sup_{[x',x_i]}f(x).
\end{align*}
定义自明。于是 $m_i=\min(m_i',m_i'')$，$M_i=\max(M_i',M_i'')$，而可见 $s_{P'}\ge s_P$，$S_{P'}\le S_P$，就是说，分划加细后下和不减而上和不增。把 $P'$ 与 $P''$ 合并成一新分划 $P$（不妨写作 $P'+P''=P$），则 $P$ 同时是 $P'$ 与 $P''$ 的加细。因此
\begin{equation}
 s_{P'}\le s_P\le S_P\le S_{P''}.
 \tag{10}
\end{equation}
就是说：不论哪一个分划的下和一定不大于任一分划的上和。于是若固定一分划 $P_0$，可见对一切分划 $P$，$s_P\le S_{P_0}$，从而 $\{s_P\}$ 有界而有上确界，记之为
\[
 \underline I=\sup_P\{s_P\},
\]
同理，$\{S_P\}$ 有下界而有下确界 $\overline I$：
\[
 \overline I=\inf_P\{S_P\}.
\]
而且
\[
 \underline I\le\overline I.
\]
$\underline I$ 和 $\overline I$ 分别称为 $f(x)$ 在 $[a,b]$ 上的下积分与上积分。它们是某集合的上、下确界，而不是序列的极限。看下面的定义即知，说积分是积分和的极限是不准确的。

\noindent\textbf{定义1}\quad 若 $\underline I=\overline I$，则称 $f(x)$ 在 $[a,b]$ 上黎曼可积（R-可积，在不产生误解时就说可积），其公共值 $I=\underline I=\overline I$ 称为 $f(x)$ 在 $[a,b]$ 上的定积分，记作
\begin{equation}
 I=\int_a^b f(x)\,dx.
 \tag{11}
\end{equation}
以上是达布（G. Darboux）对黎曼的讲法作的一个修正。黎曼本人仍是使用的积分和(7)。(7)与 $I$ 的关系以后再讲，目前我们先用达布上下和来讨论 $f(x)$ 可积的充分必要条件。

\noindent\textbf{定理1}\quad $[a,b]$ 上的有界函数 $f(x)$ 为可积的充分必要条件是：对任一 $\varepsilon>0$ 均可找到一个分划 $P$ 使
\[
 0\le S_P-s_P<\varepsilon.
\]

\noindent\textbf{证}\quad \textbf{必要性}. 设 $f(x)$ 可积，则 $I=\overline I=\underline I$。由上、下确界的定义，必可找到分划 $P'$ 和 $P''$ 使
\[
 s_{P'}>\underline I-\frac\varepsilon2=I-\frac\varepsilon2,
 \qquad
 S_{P''}<\overline I+\frac\varepsilon2=I+\frac\varepsilon2.
\]
令 $P=P'+P''$，则
\[
 0\le S_P-s_P\le S_{P''}-s_{P'}<I+\frac\varepsilon2-I+\frac\varepsilon2=\varepsilon.
\]

\textbf{充分性}. 由上下确界的定义 $\overline I\le S_P$，$\underline I\ge s_P$，但 $\overline I\ge\underline I$，所以
\[
 0\le\overline I-\underline I\le S_P-s_P<\varepsilon.
\]
由 $\varepsilon$ 之任意性有 $\overline I=\underline I$，即 $f(x)$ 可积，定理证毕。

注意到
\[
 S_P-s_P=\sum_{i=1}^{n}(M_i-m_i)(x_i-x_{i-1}).
\]
为使 $S_P-s_P$ 充分小，就应该要求在相当多的小区间 $[x_{i-1},x_i]$ 上 $M_i-m_i$ 充分小，而使 $M_i-m_i$ 不太小的小区间 $[x_{i-1},x_i]$ 之总长很小。但 $M_i-m_i$ 相当小时就很接近于函数的连续性。所以为使 $S_P-s_P$ 可以任意小，就要求在相当多的小区间上 $f(x)$ 连续，而 $f(x)$ 不连续的点不太多。这样就把 $f(x)$ 之连续性与可积性联系起来了。这一段话很不明确，读者完全可以认为是说错了。为了把这一点弄明确，我们引进函数振幅的概念。现在我们有意地设 $f(x)$ 是 $n$ 元函数。设 $f:A\subset\mathbb R^n\to\mathbb R$，$a\in A$，$f(x)$ 在 $a$ 点不连续的程度可以度量如下：令
\[
 M(f,a,\delta)=\sup\{f(x):x\in A\text{ 且 }|x-a|<\delta\},
\]
\[
 m(f,a,\delta)=\inf\{f(x):x\in A\text{ 且 }|x-a|<\delta\}.
\]
我们称
\[
 O(f,a)=\lim_{\delta\to0}\bigl[M(f,a,\delta)-m(f,a,\delta)\bigr]
\]
为 $f$ 在 $a$ 点的振幅。由于 $M(f,a,\delta)-m(f,a,\delta)$ 是非负的随 $\delta$ 不增的函数，所以 $O(f,a)$ 存在。显然 $f$ 在 $a$ 点相对于 $A$ 连续（不连续）即 $O(f,a)=0$（$\ne0$）。这里我们特意加上“相对于 $A$”几个字，因为定义 $M$ 与 $m$ 时并没有说 $x$ 在球 $|x-a|<\delta$ 内，而是在 $A\cap\{|x-a|<\delta\}$ 内（注意，我们的讨论不只对一元函数 $f(x)$ 适用，而对 $x=(x_1,\cdots,x_n)$ 时的多元函数 $f(x)$ 也适用，这时 $|x-a|<\delta$ 确实是球体而不是区间）。即只用到此球含于 $A$ 内的部分，因为 $f(x)$ 只定义于 $A$ 上，而不是定义在整个球 $|x-a|<\delta$ 上（当然也不会对此球内一切 $x$ 点 $f(x)$ 均无定义；至少对中心 $x=a$，$f(x)$ 是有定义的），这样得到的连续性是相对于 $A$ 的连续性。在一元函数情况下，若 $f(x)$ 定义在 $A=[a,b]$ 上，则在定义域的左右两端只讨论 $f(x)$ 的单侧连续性，这就是相对于 $[a,b]$ 的连续性。但在多元函数情况，由于定义域 $A$ 的构造可以十分复杂，相对于 $A$ 的连续性也复杂多了。例如，若 $a\in A$ 是一个孤立点，则 $\{x\in A,|x-a|<\delta\}$ 当 $\delta$ 充分小时就只有一个 $a$ 点，这时自然有 $O(f,a)=0$，即 $f$ 在 $a$ 点相对于 $A$ 连续。由于这个原因，在研究多元函数时，我们常限于 $f(x)$ 定义于一开集上。由于开集中的点都是内点，就不会出现上面的麻烦。总之，多元函数在边界点上的情况是要特别细心处理的。这一类问题在第六章讲到子空间拓扑时还会详细说。

\noindent\textbf{引理2}\quad 设有界函数 $f(x)$ 定义在 $[a,b]$ 上，若对一切 $x\in[a,b]$ 均有 $O(f,x)<\varepsilon$，则必存在一个分划 $P$，使
\[
 0\le S_P-s_P<2\varepsilon(b-a).
\]

\noindent\textbf{证}\quad 对一切 $x\in[a,b]$ 均存在一个开区间 $(x-\delta,x+\delta)$ 使 $M(f,x,\delta)-m(f,x,\delta)<2\varepsilon$。开区间族 $\{(x-\delta,x+\delta)\}$ 覆盖 $[a,b]$，故可由其中选出有限多个 $U_1,\cdots,U_m$ 使 $[a,b]\subset\bigcup_{k=1}^{m}U_k$。适当选分划 $P$ 可使 $P$ 之每一个小区间 $[x_{i-1},x_i]$ 均含于一个 $U_k$ 中，于是 $M_i-m_i<2\varepsilon$，而
\[
 0\le S_P-s_P=\sum_{i=1}^{n}(M_i-m_i)(x_i-x_{i-1})<2\varepsilon(b-a).
\]

下面我们就可以联系着 $f(x)$ 的连续性来讨论 $f(x)$ 可积的充分必要条件了。我们的结论是 $f(x)$ 可积的充分必要条件为其不连续点构成一个0测度集（准确些说是0勒贝格（Lebesgue）测度或0L测度集）。似乎有必要先讲什么是测度（勒贝格测度、L测度），但是事实上不需要。介绍这个概念我们又是对 $n$ 维情况来讲。

\noindent\textbf{定义2}\quad 设 $A\subset\mathbb R^n$，我们说 $A$ 具有0测度即指对任一 $\varepsilon>0$ 均可找到可数多个平行于坐标轴的闭长方体 $\{U_1,\cdots,U_m,\cdots\}$ 使 $A\subset\bigcup_{k=1}^{\infty}U_k$，而且 $\sum_{k=1}^{\infty}\operatorname{vol}(U_k)<\varepsilon$，$\operatorname{vol}(U_k)$ 即指 $U_k$ 之体积。

\noindent\textbf{注1}\quad 定义中的 $U_k$ 也可改为开长方体，只要把每个闭的 $U_k$ 代以同心、平行而稍大的开长方体即可。

\noindent\textbf{注2}\quad 任意可数点集 $\{x_k\}$ 必有0测度，因为我们可以把 $x_k$ 放在一个体积为 $\varepsilon/2^{k+1}$ 的闭长方体的中心，而这样得到的可数多个闭长方体之总体积是 $\varepsilon\sum_{k=1}^{\infty}1/2^{k+1}=\varepsilon/2<\varepsilon$，是任意小数。\jiaokan{原文此处把 $\sum_{k=1}^{\infty}2^{-(k+1)}$ 误算为1；实际为 $1/2$。结论“可数点集为0测度集”不受影响。}特别是 $[0,1]$ 中的有理数构成一个零测度集。

\noindent\textbf{注3}\quad 有限多个 $\{U_1,\cdots,U_m\}$ 也可以算作可数多个 $\{U_1,\cdots,U_m,\cdots\}$，只要假设 $U_{m+k}=\varnothing$，$k=1,2,\cdots$ 即可，但此定义中必须说可数多个 $U_k$，若只用有限多个 $U_k$ 将得到另一概念。

\noindent\textbf{定义3}\quad 设 $A\subset\mathbb R^n$，我们说 $A$ 具有0容度（0若尔当（Jordan）测度，0J测度）即是对任一 $\varepsilon>0$ 均可找到有限多个平行于坐标轴的长方体 $\{U_1,\cdots,U_m\}$ 使 $A\subset\bigcup_{k=1}^{m}U_k$，$\sum_{k=1}^{m}\operatorname{vol}(U_k)<\varepsilon$。

显然0容度集必为0测度集，但其逆不真。

如果某性质只在一个0测度集上不成立，就说它几乎处处成立，时常简写为 p.p. 成立或 a.e. 成立。

\noindent\textbf{定理3}\quad 设 $f(x)$ 是 $[a,b]$ 上的有界函数，则 $f(x)$ 在 $[a,b]$ 上可积的充分必要条件是它几乎处处相对于 $[a,b]$ 连续。

\noindent\textbf{证}\quad 先设 $f(x)$ 几乎处处连续。记 $B=\{x:x\in[a,b],f\text{ 在 }x\text{ 点不连续}\}$，所以 $B$ 有0测度。又记 $B_\varepsilon=\{x:x\in[a,b],O(f,x)\ge\varepsilon\}$，于是 $B_\varepsilon\subset B$ 也为0测度集，因此能用总长度 $<\varepsilon$ 的可数多个开的⼩区间覆盖，即是说这些小区间成为 $B_\varepsilon$ 的一个开覆盖。可是 $B_\varepsilon$ 是闭集，又是有界的：$B_\varepsilon\subset[a,b]$，所以在这个开覆盖中能取有限多个 $\{U_1,\cdots,U_m\}$ 覆盖 $B_\varepsilon$，且 $\sum_{k=1}^{m}\operatorname{vol}(U_k)<\varepsilon$。把每一个 $U_k$ 换成 $\overline U_k$ 仍可覆盖 $B_\varepsilon$，且总长度不变，所以不妨设这些 $U_k$ 均为闭区间。

对于任一个分划 $P$，有
\[
 S_P-s_P=\sum_{i=1}^{n}(M_i-m_i)(x_i-x_{i-1}).
\]
正如前面已指出的，为使 $S_P-s_P$ 充分小，我们或者要求 $M_i-m_i$ 对大多数子区间充分小，或者要求使 $M_i-m_i$ 不太小的区间总长度很小，所以我们把 $P$ 中的小区间分成两组：
\[
 S_1:\text{ 与 }B_\varepsilon\text{ 不相交的小区间},
\]
\[
 S_2:\text{ 完全含于某个 }U_k\text{ 内的小区间}.
\]
若一个小区间既不全含于任一 $U_k$ 内又与 $B_\varepsilon$ 相交，则可将它细分为两部分，一部分与 $B_\varepsilon$ 不相交，另一部分含于某个 $U_k$ 内，所以将 $P$ 的小区间分为两组是可能的。

对 $S_1$ 中的小区间，因 $O(f,x)<\varepsilon$，故由引理2，必要时对 $P$ 加细为 $P'$ 即可知，对 $P'$ 中 $S_1$ 类小区间有
\[
 \sum(M_i-m_i)(x_i-x_{i-1})<2\varepsilon(b-a).
\]
对 $S_2$ 中的小区间，因为 $f(x)$ 在 $[a,b]$ 上有界，故 $(M_i-m_i)\le2\sup_{[a,b]}|f(x)|=2M$，而对这些小区间
\[
 \sum(M_i-m_i)(x_i-x_{i-1})<2M\sum_{k=1}^{m}\operatorname{vol}(U_k)<2M\varepsilon.
\]
把这两个式子加起来即有
\[
 0\le S_{P'}-s_{P'}<2\varepsilon(b-a)+2M\varepsilon.
\]
因此，$f(x)$ 在 $[a,b]$ 上可积。

反过来设 $f(x)$ 可积，今证 $B$ 有0测度。但是 $B=B_1\cup B_{1/2}\cup\cdots\cup B_{1/m}\cup\cdots$，所以今证每个 $B_{1/m}$ 均有0测度即可。任给 $\varepsilon>0$，因已设 $f(x)$ 可积，故必有一个分划 $P$ 使 $S_P-s_P<\varepsilon/m$。用上面同样的方法将 $P$ 中的小区间分成两类，我们只看与 $B_{1/m}$ 相交的那些小区间 $[x_{i-1},x_i]$。若 $P$ 相当细，则在其上 $M_i-m_i\ge1/m$，把这些小区间合并起来，有\jiaokan{原文此处写成 $M_i-m_i\ge 2/m$，并相应写 $\frac{2}{m}\sum\operatorname{vol}(U_k)$。但 $B_{1/m}$ 的定义只给出振幅 $O(f,x)\ge1/m$；细分后能保证的下界应为 $1/m$。改正后仍足以推出 $B_{1/m}$ 为0测度集。}
\[
 \frac1m\sum\operatorname{vol}(U_k)\le\sum(M_i-m_i)(x_i-x_{i-1})=S_P-s_P<\frac\varepsilon m.
\]
所以 $\sum\operatorname{vol}(U_k)<\varepsilon$。由 $\varepsilon$ 之任意性可知 $B_{1/m}$ 为0容度的。把可数多个0容度集并起来自然可得0测度集。所以 $B$ 为0测度集而 $f(x)$ 几乎处处连续，定理证毕。

至此为止我们都是应用达布的上、下和。怎样转到黎曼所用的积分和(7)呢？

\noindent\textbf{定理4}\quad 设 $f(x)$ 在 $[a,b]$ 上有界，对任一 $\varepsilon>0$，均存在 $\delta>0$，使若某一分划 $P$ 的 $\delta_P<\delta$，就有
\[
 s_P>\underline I-\varepsilon,\qquad S_P<\overline I+\varepsilon.
\]

\noindent\textbf{证}\quad 不失一般性可设 $f(x)\ge0$（必要时用 $f(x)+C$ 代替 $f(x)$，$C$ 是充分大的正数），由 $\overline I$ 之定义，必有分划 $P_0:a=x_0<x_1<\cdots<x_n=b$，使 $S_{P_0}<\overline I+\varepsilon/2$，设 $M=\sup_{[a,b]}f(x)$，今证 $\delta=\varepsilon/[4(n+1)M]$ 即合于所求。

设 $P$ 是任一分划，和证明定理3一样，把 $P$ 的小区间分为两组：一组是完全含于 $[x_i-\delta,x_i+\delta]$ 的，其余列入第二组。于是 $S_P$ 也分成了两项 $S_P^{(1)}$ 和 $S_P^{(2)}$，$S_P^{(1)}$ 中最多有 $n+1$ 项，每一项都含了 $P_0$ 的一个分点（$P_0$ 一共有 $n+1$ 个分点），所以这些项相应于 $P$ 中的 $n+1$ 个小区间，总长为 $2(n+1)\delta=\varepsilon/(2M)$，第二组的小区间由于规定其中不能含有任一个分点，而且其长均不大于 $\delta$，所以必全位于某个小区间 $[x_{i-1},x_i]$ 内，因此
\[
 S_P=S_P^{(1)}+S_P^{(2)}\le M\cdot\frac{\varepsilon}{2M}
 +\sum_{i=1}^{n}M_i(x_i-x_{i-1})<\frac\varepsilon2+\overline I+\frac\varepsilon2=\overline I+\varepsilon.
\]
同理 $s_P>\underline I-\varepsilon$，引理证毕。

\noindent\textbf{定理5}\quad 设 $f(x)$ 在 $[a,b]$ 上有界并可积，则对任一 $\varepsilon>0$，必可找到 $\delta>0$ 使若某一分划 $P$ 的 $\delta_P<\delta$，不论如何选(7)式中的 $\xi_i$（$x_{i-1}\le\xi_i\le x_i$）都有
\[
 \left|\sum_{i=1}^{n}f(\xi_i)(x_i-x_{i-1})-\int_a^b f(x)\,dx\right|<\varepsilon.
\]
其逆亦真。

\noindent\textbf{证}\quad 因为 $m_i\le f(\xi_i)\le M_i$，所以
\[
 s_P\le\sum_{i=1}^{n}f(\xi_i)(x_i-x_{i-1})\le S_P.
\]
另一方面，由于已设 $f(x)$ 可积，故 $\underline I=\overline I=I$，而定理4又推出，只要 $\delta_P<\delta$，必有 $s_P>I-\varepsilon/2$，$S_P<I+\varepsilon/2$，所以当 $\delta_P<\delta$ 时
\[
 I-\varepsilon=\underline I-\varepsilon<s_P\le\sum_{i=1}^{n}f(\xi_i)(x_i-x_{i-1})\le S_P<\overline I+\varepsilon=I+\varepsilon,
\]
而定理得证（逆定理是明显的）。

这样我们看到，不论是用黎曼的积分和或达布的上下和来定义定积分，结果都是一样的。我们也由此看到了所谓定积分是积分和的极限这句话的确切含意。

由定理3我们就自然会提出，究竟哪一些定义于 $[a,b]$ 上的函数是可积的？当然连续函数，只有有限多个不连续点的函数，乃至有可数多个不连续点的函数都是可积的。最后这种情况是由于可数点集均为0测度集，证如下。令不连续点为 $x_1<x_2<\cdots<x_k<\cdots$，作闭区间 $U_k=[x_k-\varepsilon/2^{k+1},x_k+\varepsilon/2^{k+1}]$，显然 $\{x_1,x_2,\cdots,x_k,\cdots\}\subset\bigcup_{k=1}^{\infty}U_k$（注意 $U_k\cap U_l$ 可以是非空的），而 $\sum_{k=1}^{\infty}\operatorname{vol}(U_k)=2\sum_{k=1}^{\infty}\varepsilon/2^{k+1}=\varepsilon$。\jiaokan{原文末项误写为 $2\varepsilon$；因 $\sum_{k=1}^{\infty}2^{-(k+1)}=1/2$，故总长度为 $\varepsilon$。}证毕。具有可数多个不连续点的函数的一个重要情况是单调函数 $f(x)$。这是因为 $\inf_{[a,b]}f(x)=f(a)$，$\sup_{[a,b]}f(x)=f(b)$（设 $f(x)$ 为不减的）。记 $M=f(b)-f(a)$，$B=\{x:f\text{ 在 }x\text{ 点不连续}\}$，因为单调不减函数的每个不连续点 $x$ 都是第一类间断点，而 $O(f,x)=f(x+0)-f(x-0)$。依照前面的记号 $B=\bigcup_{m=1}^{\infty}B_m$，$B_m=\{x:O(f,x)\ge M/m\}$，则 $B_m$ 必为有限集，因为不然，$B_m$ 中至少有可数多个点 $x_1<x_2<\cdots$，而
\[
 f(b)-f(a)\ge\sum_{i=1}^{\infty}[f(x_i+0)-f(x_i-0)]\ge\sum_{i=1}^{\infty}\frac Mm=+\infty.
\]
$B$ 作为可数多个有限集之并当然最多是可数的。

\noindent\textbf{推论6}\quad $[a,b]$ 上的单调函数必为可积的。

\noindent\textbf{3. 微积分的基本定理}\quad
以上我们系统地讨论了函数可积性问题，下面转向另一个问题：积分和微分的关系问题。人们心目中都以为二者关系很简单，即二者互为逆运算；认识到二者互为逆运算十分重要，因为微分表反过来看就是积分表。自此以后，人们就知道积分是怎样计算的。如果没有这样一个表，遇到任何一个问题都得利用第一章里介绍的希腊人的穷竭法去处理，微积分就决不可能成为研究种种实际问题的工具，不可能为物理学家与工程师所掌握。所以牛顿和莱布尼茨把他们“发现”的这个公式
\begin{equation}
 \int_a^b F'(x)\,dx=F(b)-F(a)
 \tag{12}
\end{equation}
称为微积分的基本定理。但是在牛顿和莱布尼茨的时代，对数学分析的基本概念还很不清楚，所以他们的论证也很不清楚。直到柯西、黎曼的时代，这个定理才有明确的证明。柯西对连续的 $F'(x)$ 证明了(12)，达布则对一般可积的 $F'(x)$ 证明了它。

我们首先直观地分析一下(12)的含意。因为微分和积分是如此不同的概念，而有如(12)这样的公式将它们联结起来当然是很有意思的。为此，我们先把“极限”的成分排除掉，而把“连续”的东西代以相应的“离散”的东西如下：
\[
 \int_a^b\quad\text{代以}\quad\sum_{i=1}^{n},
\]
\[
 f'(x)\quad\text{代以}\quad\left(\frac{\Delta f}{\Delta x}\right)_i
 =\frac{f(x_i)-f(x_{i-1})}{x_i-x_{i-1}},
\]
\[
 dx\quad\text{代以}\quad(\Delta x)_i=x_i-x_{i-1}.
\]
这样一来，(12)就化为一个几乎自明的公式
\[
 \sum_{i=1}^{n}\frac{\Delta f}{\Delta x_i}(\Delta x)_i
 =\sum_{i=1}^{n}(\Delta f)_i
 =\sum_{i=1}^{n}[f(x_i)-f(x_{i-1})]
 =f(x_n)-f(x_0)=f(b)-f(a).
\]
可见，(12)原来是一个很简单的组合学公式。组合学成分是积分的重要因素，在本书最后一章我们还要讨论这一件事。但现在我们把极限成分放进去，这时就会发现(12)中其实包含了两个不同的问题。

（1）若 $f(x)$ 可积，则它的不定积分 $\int_a^x f(\xi)\,d\xi$ 是否 $f(x)$ 的原函数？$f(x)$ 的原函数这里就是指处处（而不仅是“几乎处处”）适合关系式 $F'(x)=f(x)$ 的函数 $F(x)$。

（2）若 $f(x)$ 有原函数 $F(x)$，$f(x)$ 是否可积？

可惜，两个问题的答案均是否定的，为此，我们看两个反例：

（1）$f(x)$ 可积但没有原函数的例子。在区间 $[0,1]$ 中考虑一个可数的处处稠密的点集 $\{a_k\}$，例如有理点的集合即是一个可数的而且处处稠密的点集。再看一个收敛的正项级数 $\sum_{k=1}^{\infty}p_k=p<+\infty$。我们定义一个函数 $f(x)$ 如下：
\begin{equation}
 f(x)=\sum_{x\ge a_{i_k}}p_{i_k}.
 \tag{13}
\end{equation}
就是说，看 $x$ 之左方（包括 $x$）：$[0,x]$ 包含了哪些 $a_{i_k}$，若 $a_{i_k}\in[0,x]$，则将级数中有相同标号的项 $p_{i_k}$ 列入(13)右方。因为 $\sum_{k=1}^{\infty}p_k=p<\infty$ 是正项级数，故其部分项（尽管是无穷多项）所成的级数(13)仍收敛，从而(13)定义的函数 $f(x)$ 是有意义的。不但如此，当 $x$ 增加时，适合不等式 $a_{i_k}\le x$ 的 $a_{i_k}$ 更多，因此(13)中还要多加几个项：
\begin{equation}
 f(x+h)-f(x)=\sum_{x+h\ge a_{i_k}}p_{i_k}-\sum_{x\ge a_{i_k}}p_{i_k}
 =\sum_{x+h\ge a_{i_k}>x}p_{i_k}\ge0.
 \tag{14}
\end{equation}
因此，$f(x)$ 不减而为单调函数，从而 $f(x)$ 可积。

但是 $f(x)$ 不可能有原函数，因为若 $f(x)$ 有原函数 $F(x)$，当有 $F'(x)=f(x)$。对于(13)中的 $f(x)$，每个 $a_k$ 都是跳跃间断点；而导函数具有介值性质（Darboux 性质），不可能具有这样的跳跃间断。因此不存在处处满足 $F'(x)=f(x)$ 的原函数 $F(x)$。也就是说，若试图构造这样的 $F$，它至少在稠密点集 $\{a_k\}$ 上不能具有所要求的导数。\jiaokan{原文此处写“而 $F(x)$ 在 $a_k$ 处没有 $F'(x)$。这与 $F(x)$ 是 $f(x)$ 的原函数矛盾。$F(x)$ 不但不是几乎处处可求导，而且在一个稠密集上不可导”。仅由 $f$ 在 $a_k$ 有跳跃不能直接推出 $F$ 在该点不可导；严格理由是导函数具有 Darboux 介值性质，而(13)的 $f$ 在每个 $a_k$ 有跳跃。另因 $\{a_k\}$ 可数，原文“不是几乎处处可求导”的表述也不正确。}

（2）$f(x)$ 有原函数但本身不可积的例子。令
\begin{equation}
 F(x)=
 \begin{cases}
 x^2\sin\dfrac1{x^2},&x\ne0,\\[4pt]
 0,&x=0.
 \end{cases}
 \tag{15}
\end{equation}
令 $f(x)=F'(x)$，可见 $f(x)$ 有原函数，但 $f(x)$ 在 $x=0$ 处是无界的，因此不是可积的。

联系着这两个例子来看微积分的基本定理则应该表述如下：

\noindent\textbf{定理7（微积分的基本定理，达布）}\quad 设

（1）$f(x)$ 在 $(A,B)$ 上可积；

（2）$f(x)$ 在 $(A,B)$ 上有原函数 $F(x)$，

则对任一闭区间 $[a,b]\subset(A,B)$ 有
\begin{equation}
 F(b)-F(a)=\int_a^b f(x)\,dx.
 \tag{16}
\end{equation}

\noindent\textbf{证}\quad 作任一分划 $P:a=x_0<x_1<\cdots<x_n=b$，有
\[
 F(b)-F(a)=\sum_{i=1}^{n}[F(x_i)-F(x_{i-1})]
 =\sum_{i=1}^{n}F'(\xi_i)(x_i-x_{i-1}),
 \qquad x_{i-1}<\xi_i<x_i.
\]
此式右方是 $F'(x)=f(x)$ 之积分和，故当 $\delta_P\to0$ 时，其极限为 $\int_a^b f(x)\,dx$。但式左 $F(b)-F(a)$ 之值不随分划的选择而变，从而(16)成立，证毕。
